\section{Fundamentos de movilidad, percepción y parámetros}
\label{sec:fundamentos}

La \cref{sec:dinamica} introdujo la dinámica de rueda y de carga, y
\cref{sec:movimiento} dos lemas sobre orientación y curvatura. Falta la base
que explica \emph{por qué} esas restricciones tienen la forma que tienen, y
falta declarar dos supuestos que el resto del trabajo usa sin nombrarlos: que la
pose se conoce y que los parámetros de la carga se conocen.

\subsection{Clasificación de movilidad}

La estructura cinemática de un robot con ruedas se caracteriza por dos enteros.
Sea $C_1^{f}(\beta_s)$ la matriz que apila las restricciones de deslizamiento de
las ruedas fijas y direccionables.

\begin{definicion}[Grados de movilidad, direccionabilidad y maniobrabilidad]
\label{def:movilidad}
\[
  \delta_m=\dim\mathcal N\bigl(C_1^{f}(\beta_s)\bigr),\qquad
  \delta_s=\operatorname{rank}\bigl(C_1^{s}(\beta_s)\bigr),\qquad
  \delta_M=\delta_m+\delta_s ,
\]
donde $\delta_m$ cuenta los grados de libertad del chasis manipulables
\emph{de inmediato} mediante velocidades de rueda, y $\delta_s$ los parámetros
de dirección independientes \cite{siegwart2011introduction}.
\end{definicion}

Para el AMR de tracción diferencial, las dos ruedas fijas comparten eje
horizontal y la segunda no añade restricción independiente: $\delta_m=2$,
$\delta_s=0$, luego $\delta_M=2$. El robot controla directamente su velocidad
longitudinal y su velocidad de giro, y su centro instantáneo de rotación está
obligado a caer sobre la recta que prolonga el eje de las ruedas.

\begin{observacion}[De dónde vienen los lemas de movimiento]
\label{obs:origen-lemas}
$\delta_m=2<3$ es exactamente el motivo de \cref{le:orientacion}: con dos grados
manipulables sobre un chasis de tres, la dirección de empuje no es libre y hay
que orientarse. Y $\delta_s=0$ es el motivo de \cref{le:curvatura}: sin ruedas
direccionables, curvatura y velocidad se acoplan a través de las mismas dos
entradas. Un AMR omnidireccional tendría $\delta_m=3$ y ninguno de los dos lemas
sería necesario; el precio es mecánico, no algorítmico.
\end{observacion}

\begin{observacion}[Consecuencia para la matriz de agarre]
\label{obs:movilidad-agarre}
La columna $g_{ks}$ de \eqref{eq:columna} supone que el puesto entrega su
esfuerzo en la dirección $n_{ks}$. Con $\delta_M=2$ esa dirección no es
instantáneamente arbitraria: pertenece al cono admisible del chasis en su
orientación actual. Por tanto el conjunto realizable de \cref{th:span} debe
entenderse sobre las direcciones \emph{alcanzables tras reorientación}, y el
tiempo de esa reorientación entra en el presupuesto de la maniobra. La
caracterización por span positivo se conserva; lo que cambia es que sus columnas
son función de las poses, no constantes.
\end{observacion}

\subsection{El supuesto de pose conocida}

Todo el aparato de barreras evalúa $h_{ij}=\|p_i-p_j\|^2-d_{\min}^2$ sobre poses
\emph{estimadas}. La estimación de pose en interiores se resuelve con filtros
bayesianos —Markov, Kalman extendido, partículas— sobre el mismo tipo de medidas
telemétricas que alimentan \cref{th:mapa} \cite{siegwart2011introduction,thrun2005probabilistic}.

\begin{supuesto}[Pose con error acotado]
\label{s:pose}
Existe $\varepsilon_p>0$ tal que $\|\hat p_i-p_i\|\le\varepsilon_p$ para todo $i$
y todo $t$.
\end{supuesto}

\begin{proposicion}[Estrechamiento de la barrera por error de pose]
\label{pr:pose}
Bajo \cref{s:pose}, imponer la barrera estrechada
$\hat h_{ij}=\|\hat p_i-\hat p_j\|^2-(d_{\min}+2\varepsilon_p)^2\ge0$ sobre las
poses estimadas garantiza $h_{ij}\ge0$ sobre las poses reales.
\end{proposicion}

\begin{proof}
Por desigualdad triangular,
$\|p_i-p_j\|\ge\|\hat p_i-\hat p_j\|-2\varepsilon_p$. Si
$\|\hat p_i-\hat p_j\|\ge d_{\min}+2\varepsilon_p$, entonces
$\|p_i-p_j\|\ge d_{\min}$.
\end{proof}

\begin{observacion}[El coste se acumula]
\label{obs:coste-pose}
\Cref{pr:pose} se compone con \cref{le:muestras}: el margen total que hay que
reservar es $2\varepsilon_p$ por localización más $L_\Psi\Delta t$ por muestreo.
Y compite con \cref{pr:compatibilidad}: cada metro de estrechamiento por
seguridad reduce el anillo admisible frente a la conectividad, con lo que la
cota de grado del árbol generador empeora. La precisión de localización no es un
detalle de implementación: entra en la cota de grado de la topología.
\end{observacion}

\subsection{Parámetros desconocidos y equivalencia cierta}

\Cref{pr:monotonia-info} trata la incertidumbre paramétrica como conjunto
$\Theta$. La alternativa habitual en control es estimar un valor nominal
$\hat\theta$ y actuar como si fuera exacto. Conviene separar ambos objetos con
precisión, porque hacen cosas distintas.

Para un error de seguimiento $s$ cuya dinámica admite la forma regresora
$M(q)\dot s+C(q,\dot q)s=Y(q,\dot q,v_r,\dot v_r)\tilde\theta-Ks+r$
con $\tilde\theta=\hat\theta-\theta$, la ley $\dot{\hat\theta}=-\Gamma Y^{\!\top}s$
y la antisimetría de $\dot M-2C$ dan
\begin{equation}
  V=\tfrac12 s^{\!\top}Ms+\tfrac12\tilde\theta^{\!\top}\Gamma^{-1}\tilde\theta,
  \qquad
  \dot V=-s^{\!\top}Ks+s^{\!\top}r .
  \label{eq:adaptativo}
\end{equation}

\begin{observacion}[Qué da y qué no da la adaptación]
\label{obs:adaptacion}
Con $r=0$, \eqref{eq:adaptativo} da disipación, y con
$k=\lambda_{\min}(K)$ se obtiene
$\dot V\le-\tfrac{k}{2}\|s\|^2+\tfrac{1}{2k}\|r\|^2$. Eso \emph{no} demuestra
estabilidad entrada--estado del vector completo de errores, porque el término
negativo no domina el error paramétrico $\tilde\theta$. Dos consecuencias.
Primera: saturación y pérdida de contacto entran en $r$ y no desaparecen por
estimar mejor la masa; el certificado que las cubre es
\cref{th:autoridad}, no \eqref{eq:adaptativo}. Segunda: seguir bien no implica
identificar (\cref{obs:identificacion}), de modo que $\hat\theta$ sirve para
rendimiento y $\Theta$ para garantías. Confundirlos produce certificados que
descansan sobre una estimación puntual sin cobertura.
\end{observacion}

\begin{observacion}[Qué no aporta la biblioteca consultada]
\label{obs:biblioteca-fisica}
De las obras de robótica móvil y control cooperativo revisadas, ninguna formula
el reparto de esfuerzo mediante el \emph{wrench} resultante sobre una carga
común: tratan formación, seguimiento y consenso sobre poses, no asignación de
esfuerzos con unilateralidad y conos de fricción. La capa de contacto de este
trabajo no tiene precedente directo en esa literatura, y esa es la frontera
concreta que separa el control de formación del transporte cooperativo.
\end{observacion}

\begin{observacion}[Un equilibrio generalizado puede ser socialmente inútil]
\label{obs:gne-ineficiente}
La definición de equilibrio de Nash generalizado se vacía cuando la
restricción compartida bloquea toda desviación. Con $x_1+x_2=1$,
$0\le x_i\le1$, $J_1=x_1^2/2$ y $J_2=3x_2^2/2$, fijado el otro jugador la
igualdad deja una sola acción admisible, luego \emph{todo} punto factible es
un equilibrio generalizado; el único equilibrio variacional es $(3/4,1/4)$,
con coste social $3/8$ frente a $1/2$ en el equilibrio generalizado $(1,0)$
\cite{mayorga2026compendio}. La selección variacional no es un refinamiento
técnico: es lo que hace que el equilibrio con multiplicador común de
\cref{th:potencial} tenga contenido cuando la restricción compartida es
activa.
\end{observacion}

\begin{observacion}[El pago vectorial importa más que el protocolo]
\label{obs:pago-vectorial}
Sobre ciento veinte mundos de una relajación continua de asignación, tres
inicializaciones y tres mil quinientas iteraciones, las tres dinámicas
—replicador, Smith y BNN— dan brechas medianas respecto de la referencia
convexa de $1{,}211$, $1{,}211$ y $1{,}211$ con pago escalar, y de
$1{,}2\times10^{-4}$, $3{,}9\times10^{-4}$ y $7{,}0\times10^{-3}$ con pago
vectorial \cite{mayorga2026compendio}. Representar correctamente la
contribución vectorial de cada AMR a cada carga cambia el resultado en cuatro
órdenes de magnitud; cambiar de protocolo con pago escalar no lo cambia. Es
la versión empírica de \cref{th:noequiv}: la heterogeneidad no se resume en
un escalar. Y a presupuesto finito ninguna alcanzó la brecha variacional de
$10^{-4}$, de modo que las tres difieren en velocidad aunque coincidan en
destino.
\end{observacion}
